\documentclass[11pt]{article}

\usepackage[a4paper,margin=1in]{geometry}
\usepackage{amsmath,amssymb,amsthm,mathtools}
\usepackage{graphicx}
\usepackage[hidelinks]{hyperref}

\newtheorem{theorem}{Theorem}
\newtheorem{lemma}{Lemma}
\newtheorem{corollary}{Corollary}
\theoremstyle{definition}
\newtheorem{definition}{Definition}
\newtheorem{remark}{Remark}

\newcommand{\R}{\mathbb R}
\newcommand{\Ffin}{\mathcal F}
\newcommand{\norm}[1]{\left\lVert #1\right\rVert_\infty}
\newcommand{\supp}{\operatorname{supp}}
\newcommand{\sgn}{\operatorname{sgn}}
\newcommand{\row}{\operatorname{row}}
\newcommand{\Span}{\operatorname{span}}

\title{A Complete Characterization of Property-\((R_1)\) Kernels in Finite-Dimensional \(\ell_\infty\)}
\author{Automatic math-discovery run \texttt{fa\_banach\_001}, lane 14}
\date{August 9, 2026}

\begin{document}
\maketitle

\section*{Result Status}

\textbf{Candidate full solution, likely valid, pending expert review.}
The result answers the characterization problem in Remark 5.2(a) of
Das--Paul \cite{DasPaul}.  The answer depends only on the row space of the
given functionals: after invertible row operations, their matrix must be a
signed coordinate-incidence matrix.

\section{Source Problem}

The source is Syamantak Das and Tanmoy Paul,
\emph{On property-\((R_1)\) and relative Chebyshev centers in Banach
spaces-II}, arXiv:2207.09623 \cite{DasPaul}.  The question appears in Remark
5.2(a), page 17 of the arXiv PDF.

\begin{center}
\includegraphics[width=\linewidth]{figures/open_problem_crop.png}
\end{center}

In the finite-dimensional part of that remark, the problem is to characterize
families
\[
 \alpha_1,\ldots,\alpha_m\in\ell_1^n
\]
satisfying
\[
 \dim\left(\ell_\infty^n\Big/\bigcap_i\ker\alpha_i\right)=m,
\]
for which
\[
 V:=\bigcap_{i=1}^m\ker\alpha_i
\]
has property \((R_1)\) in \(\ell_\infty^n\), with respect to all finite
subsets.

\section{Definitions}

For a nonempty finite set \(E\subset\ell_\infty^n\), put
\[
 S_s(E):=\bigcap_{x\in E}B_\infty(x,s),
 \qquad r(v,E):=\max_{x\in E}\norm{v-x}.
\]
The triplet \((\ell_\infty^n,V,\Ffin(\ell_\infty^n))\) has
\emph{strong property \((R_1)\)} if, whenever \(v\in V\), \(r_1,r_2>0\),
and \(E\) is finite,
\[
 r(v,E)\le r_1+r_2,\qquad S_{r_2}(E)\cap V\ne\varnothing,
\]
one has
\[
 V\cap B_\infty(v,r_1)\cap S_{r_2}(E)\ne\varnothing.
\]
Ordinary property \((R_1)\) is the corresponding approximate/open-radius
condition stated in the source paper.

For \(r\ge0\), define coordinatewise clipping by
\[
 (T_r x)_j:=\sgn(x_j)\min\{|x_j|,r\}.
\]
We identify \((\ell_\infty^n)^*=\ell_1^n\) with \(\R^n\), using the
standard coordinate vectors \(e_1,\ldots,e_n\).

\section{Proof Intuition}

Intersections of equal-radius \(\ell_\infty\)-balls are precisely bounded
axis-parallel boxes.  Property \((R_1)\) says that if such a box meets the
subspace and also meets a cube centered at a point of the subspace, then all
three sets meet.  Centering at zero turns this into a nearest-point statement
for boxes.  The closest point of a box to zero is obtained by clipping any
point in the box toward the coordinate cube.  Thus the entire metric condition
is equivalent to the algebraic rule \(T_r(V)\subset V\).

A truncation-stable linear subspace is rigid.  A nonzero vector with smallest
support clips to a sign vector on that support.  If another vector were not a
scalar multiple of that sign pattern on the same support, clipping a small
perturbation would create a nonzero vector with strictly smaller support.
Removing the resulting signed diagonal block and iterating gives the full
decomposition.  Taking annihilators converts the decomposition into the
promised condition on the \(\alpha_i\).

\section{Main Theorem}

\begin{theorem}[Complete characterization]\label{thm:main}
Let \(\alpha_1,\ldots,\alpha_m\in\ell_1^n\) be linearly independent, let
\(A\) be the \(m\times n\) matrix with these vectors as rows, and let
\[
 V=\ker A=\bigcap_{i=1}^m\ker\alpha_i\subset\ell_\infty^n.
\]
The following are equivalent.
\begin{enumerate}
\item \((\ell_\infty^n,V,\Ffin(\ell_\infty^n))\) has property \((R_1)\).
\item The same triplet has strong property \((R_1)\).
\item \(T_r(V)\subset V\) for every \(r\ge0\).
\item After a permutation of coordinates, there are pairwise disjoint blocks
      \(B_1,\ldots,B_q\), sign vectors
      \(s^{(a)}\in\{\pm1\}^{B_a}\), a set \(Z\) of zero coordinates, and
      a set \(J\) of free singleton coordinates such that
      \[
      V=\{x\in\R^n:x_j=0\ (j\in Z),\quad
      x|_{B_a}=t_a s^{(a)}\ (t_a\in\R),\quad
      x_j\text{ arbitrary }(j\in J)\}.
      \]
\item The row space \(\row(A)=V^\perp\) admits a basis whose members, after
      individual nonzero rescaling, have one of the forms
      \[
      e_j,\qquad e_j+e_k,\qquad e_j-e_k\quad(j\ne k).
      \]
\end{enumerate}
Consequently, condition (5) is the requested characterization of the family
\(\alpha_1,\ldots,\alpha_m\).  Equivalently, invertible row operations turn
\(A\) into a matrix whose rows have either one nonzero entry or two entries of
equal absolute value.
\end{theorem}

The proof is divided into three elementary lemmas.

\section{From Property \texorpdfstring{\((R_1)\)}{(R1)} to Boxes}

\begin{lemma}\label{lem:strong}
For a linear subspace \(V\subset\ell_\infty^n\), property \((R_1)\) for all
finite subsets is equivalent to strong property \((R_1)\).
\end{lemma}

\begin{proof}
The closed unit ball \(B_V\) is compact.  For every finite \(E\), the function
\(y\mapsto r(y,E)\) is continuous, so it attains its minimum on \(B_V\).
Thus \(\operatorname{Cent}_{B_V}(E)\ne\varnothing\).  The equivalence stated
in \cite[Theorem 1.4]{DasPaul} says that strong property \((R_1)\) is
equivalent to ordinary property \((R_1)\) together with this center-existence
condition.  The conclusion follows.
\end{proof}

Call a product \(Q=\prod_{j=1}^n[a_j,b_j]\), with \(a_j\le b_j\), a bounded
coordinate box.

\begin{lemma}[Box reformulation]\label{lem:boxes}
Strong property \((R_1)\) is equivalent to the following condition:
for every bounded coordinate box \(Q\), every \(v\in V\), and every \(r>0\),
\[
 Q\cap V\ne\varnothing,\quad Q\cap B_\infty(v,r)\ne\varnothing
 \quad\Longrightarrow\quad
 Q\cap V\cap B_\infty(v,r)\ne\varnothing.
\tag{B}
\]
\end{lemma}

\begin{proof}
First note that every \(S_s(E)\) is a bounded coordinate box.  Conversely,
let \(Q=\prod_j[a_j,b_j]\).  Choose
\[
 s>0,\qquad 2s\ge\max_j(b_j-a_j),
 \qquad c_j=\frac{a_j+b_j}{2}.
\]
For each \(j\), form two points \(p^{j,+},p^{j,-}\) by taking midpoint
coordinates except at \(j\), and setting
\[
 p^{j,+}_j=a_j+s,\qquad p^{j,-}_j=b_j-s.
\]
For \(E=\{p^{j,+},p^{j,-}:1\le j\le n\}\), direct coordinatewise
calculation gives \(S_s(E)=Q\).

Suppose strong property \((R_1)\) holds, and let
\(z\in Q\cap B_\infty(v,r)\).  With the representation \(Q=S_s(E)\), for
every \(x\in E\) we have
\[
 \norm{v-x}\le\norm{v-z}+\norm{z-x}\le r+s.
\]
Together with \(Q\cap V\ne\varnothing\), strong property \((R_1)\) gives
(B).

Conversely, suppose (B) holds and the hypotheses of strong property
\((R_1)\) are satisfied.  Put \(Q=S_{r_2}(E)\).  In coordinate \(j\),
\[
 Q_j=\left[\max_{x\in E}(x_j-r_2),\ 
              \min_{x\in E}(x_j+r_2)\right].
\]
The inequalities \(\norm{v-x}\le r_1+r_2\), for all \(x\in E\), say exactly
that \(Q_j\) intersects \([v_j-r_1,v_j+r_1]\) for every \(j\).  Hence
\(Q\cap B_\infty(v,r_1)\ne\varnothing\).  Applying (B) proves strong
property \((R_1)\).
\end{proof}

\section{Boxes Are Coordinatewise Clipping}

\begin{lemma}[Clipping criterion]\label{lem:clip}
Condition \emph{(B)} is equivalent to \(T_r(V)\subset V\) for every \(r\ge0\).
\end{lemma}

\begin{proof}
Because \(V\) is linear, translating \(Q\) by \(-v\) shows that it is enough
to use (B) with \(v=0\).

Assume first that \(T_r(V)\subset V\).  Let \(Q\cap V\ne\varnothing\), choose
\(x\in Q\cap V\), and suppose \(Q\cap B_\infty(0,r)\ne\varnothing\).  For
each coordinate interval \([a_j,b_j]\) of \(Q\), the latter condition says
\([a_j,b_j]\cap[-r,r]\ne\varnothing\).  Moving \(x_j\) toward
\([-r,r]\) therefore stays inside \([a_j,b_j]\), so \(T_r x\in Q\).
By assumption \(T_r x\in V\), and by definition \(T_r x\in B_\infty(0,r)\).
Thus (B) holds.

Conversely, fix \(x\in V\) and \(r\ge0\).  The cases \(r=0\) and
\(r\ge\norm{x}\) are immediate.  Otherwise define the box
\(Q(x,r)=\prod_j I_j\), where
\[
 I_j=
 \begin{cases}
 [r,x_j],&x_j>r,\\
 [x_j,-r],&x_j<-r,\\
 \{x_j\},&|x_j|\le r.
 \end{cases}
\]
It contains \(x\in V\), while
\[
 Q(x,r)\cap B_\infty(0,r)=\{T_r x\}.
\]
Condition (B) forces this unique point into \(V\).  Hence \(T_r(V)\subset V\).
\end{proof}

\section{Classification of Clipping-Stable Subspaces}

\begin{lemma}[Minimal-support decomposition]\label{lem:structure}
A linear subspace \(V\subset\R^n\) is invariant under every \(T_r\) if and
only if it has the signed-block decomposition in Theorem \ref{thm:main}(4).
\end{lemma}

\begin{proof}
The stated block spaces are visibly clipping-stable: clipping a signed block
\(t s\) gives \(\sgn(t)\min\{|t|,r\}s\), and clipping preserves zero and
singleton blocks.

For the converse, suppose \(V\ne\{0\}\), and choose \(x\in V\setminus\{0\}\)
for which \(|\supp x|\) is minimal.  If
\(0<r<\min_{j\in\supp x}|x_j|\), then
\[
 s:=r^{-1}T_r x\in V
\]
is a sign vector supported on the same set \(B:=\supp x\).

We claim that for every \(y\in V\) there is a scalar \(c(y)\) such that
\[
 y_j=c(y)s_j\qquad(j\in B).
\tag{1}
\]
Indeed, suppose the numbers \(s_jy_j\), \(j\in B\), are not all equal.
Choose \(t>0\) so small that every coordinate of \(s+ty\) on \(B\) keeps
the sign \(s_j\), and every coordinate off \(B\) has absolute value smaller
than every coordinate on \(B\).  On \(B\),
\[
 |s_j+ty_j|=1+t s_jy_j.
\]
Choose a clipping level \(\rho\) strictly between two distinct values in this
finite list, avoiding all of them.  Then
\[
 w:=T_\rho(s+ty)-(s+ty)\in V
\]
is nonzero, has no support off \(B\), and is supported on the nonempty proper
subset of indices where \(|s_j+ty_j|>\rho\).  This contradicts the minimality
of \(B\), proving (1).

It follows that
\[
 V=\Span\{s\}\oplus V_0,
 \qquad V_0:=\{y\in V:y|_B=0\}.
\]
The second summand is again invariant under all coordinatewise clippings and
is supported off \(B\).  Iteration terminates after finitely many steps and
gives disjoint signed diagonal blocks.  Coordinates never appearing in a
support are zero on \(V\); blocks of size one are precisely free coordinate
directions.  This is the required decomposition.
\end{proof}

\begin{proof}[Proof of Theorem \ref{thm:main}]
Lemmas \ref{lem:strong}, \ref{lem:boxes}, \ref{lem:clip}, and
\ref{lem:structure} prove the equivalence of (1)--(4).

It remains to translate (4) into (5).  A zero coordinate \(j\in Z\) is cut
out by the normal \(e_j\).  For each signed block \(B_a\), choose a base
index \(j_a\in B_a\).  The block condition
\(x|_{B_a}=t_a s^{(a)}\) is cut out by the independent normals
\[
 s^{(a)}_{j_a}e_{j_a}-s^{(a)}_j e_j,
 \qquad j\in B_a\setminus\{j_a\}.
\]
Thus \(V^\perp\) has a basis of the form in (5).

Conversely, suppose \(V^\perp\) is spanned by normals \(e_j\) and
\(e_j\pm e_k\).  Regard every two-coordinate normal as a signed edge between
vertices \(j,k\).  In each connected component the equations determine every
coordinate as a fixed sign times one scalar, unless a one-coordinate equation
or an inconsistent signed cycle pins the entire component to zero.  Isolated
unconstrained vertices are free singleton blocks.  Hence \(V\) has the
decomposition in (4).

Finally, \(\row(A)=V^\perp\), and its dimension is \(m\).  A spanning family
of elementary normals therefore contains an elementary-normal basis, while
any such basis is obtained from the original rows of \(A\) by invertible row
operations.  This proves (5) and the final formulation.
\end{proof}

\begin{corollary}[Concrete obstruction]
The line \(V=\ker(1,2)=\Span\{(2,-1)\}\subset\ell_\infty^2\) does not have
property \((R_1)\).
\end{corollary}

\begin{proof}
For \(x=(2,-1)\in V\), one has \(T_1x=(1,-1)\notin V\).  Equivalently, the
box \(Q=[1,2]\times\{-1\}\) meets \(V\) at \(x\), but its unique point in
the unit ball is \((1,-1)\notin V\).  Lemma \ref{lem:clip} applies.
\end{proof}

\section{Relation to Earlier Literature}

Miesch--Pav\'on \cite[Theorem 4.7]{MieschPavon} characterize weakly
externally hyperconvex linear subspaces of \(\ell_\infty^n\).  Their answer is
the same signed-coordinate class as Theorem \ref{thm:main}(4)--(5): precisely
the intersections of hyperplanes with normals \(e_j\) or \(e_j\pm e_k\).
Their paper does not discuss property \((R_1)\), and their theorem is not used
in the proof above.  The coincidence supplies an independent structural
cross-check; the box/clipping equivalence identifies property \((R_1)\) with
that previously occurring class in finite-dimensional \(\ell_\infty\).

\section{Verification Notes}

The proof was audited around four possible failure points.
\begin{enumerate}
\item Ordinary property becomes strong property because \(B_V\) is compact;
      this is exactly the extra hypothesis in \cite[Theorem 1.4]{DasPaul}.
\item The common-radius representation really includes every bounded box,
      including degenerate intervals, because \(b_j-s\le a_j+s\).
\item In Lemma \ref{lem:structure}, \(t\) is chosen before \(\rho\), so no
      coordinate outside the minimal support is clipped.  The clipping
      difference is therefore supported on a strict nonempty subset of that
      support.
\item The criterion must concern \(\row(A)\), not each originally presented
      \(\alpha_i\).  Property \((R_1)\) depends only on \(\ker A\), and
      invertible row operations do not change that kernel.
\end{enumerate}
No numerical or computer-assisted step is used in the proof.  The example
\(\ker(1,2)\) is a direct sanity check showing how a forbidden coefficient
ratio produces a failed clipping.

\section{Novelty Check and Limitations}

Before promotion, the run's registry, solution, attempt, and proof-gap indexes
were searched for arXiv:2207.09623, the paper title, property \((R_1)\), and
the finite-dimensional \(\ell_\infty\) kernel problem; no duplicate result was
found.  Targeted web searches on August 9, 2026 used the exact source sentence
and close combinations of ``property \((R_1)\)'', ``finite dimensional'',
``\(\ell_\infty\)'', ``kernel'', and ``characterization''.  They found the
source paper and general Chebyshev-center literature, but no later paper
stating this characterization.

The closest primary result found was the earlier Miesch--Pav\'on theorem just
discussed.  Because it concerns weak external hyperconvexity rather than
property \((R_1)\), it is recorded as related literature rather than an
existing answer.  Novelty confidence is therefore moderate, not definitive:
an expert should check older terminology around the \(1\tfrac{1}{2}\)-ball property
and ball-intersection subspaces for an implicit equivalent formulation.

The theorem is real and finite-dimensional, exactly as in the source remark.
No claim is made here for complex scalars or arbitrary subspaces of infinite
\(\ell_\infty\).

\section{Human Review Recommendation}

Send for expert review as a candidate full solution.  The most useful review
is to verify the equivalence between the source's ordinary and strong
definitions, then inspect the minimal-support clipping argument.  A secondary
literature reviewer should decide whether the new box/clipping identification
is already implicit in the older \(1\tfrac{1}{2}\)-ball literature; this affects
novelty classification but not the self-contained proof.

\begin{thebibliography}{9}

\bibitem{DasPaul}
S.~Das and T.~Paul,
\emph{On property-\((R_1)\) and relative Chebyshev centers in Banach
spaces-II},
arXiv:2207.09623; Quaestiones Mathematicae \textbf{47} (2024), 461--476.

\bibitem{MieschPavon}
B.~Miesch and M.~Pav\'on,
\emph{Weakly externally hyperconvex subsets and hyperconvex gluings},
arXiv:1507.07795, Theorem 4.7.

\end{thebibliography}

\end{document}
